\subsection{从军粮到衣料的生产链}
军队的供给不只有粮食。布帛、皮革、铁器、箭矢、马具和修理所需的材料，都要经过生产、征取、运输和分发。官方材料较容易记录某次调发或某项禁令，较少记录一匹布从家庭劳作到仓场之间的具体步骤。军事史若只写将领和战役，就会把这些生产链从战争中抹去；财政史若只写税额，也会忽略物品在道路和仓储中可能损耗、折换或被截留。

手工业者和运输者在这种链条中的位置通常不清晰。城市作坊、乡村生产、军营修造和市场交换可能互相连接，却不能由某一处遗址直接推出全国性的生产组织。考古材料可以让窑址、冶铸或仓储遗存出现，必须结合断代、层位和发表情况谨慎使用。它们最有力的作用，是提醒我们物资并不是从“地方”自动流入“国家”，而要经过掌握技能、道路和文书的人。

晚唐各政权和军镇争夺的因此不只是土地或城市，也包括使这些生产链可被征取的关口。一个节度使控制某处仓场、一段运河或一座市场，可能获得军粮与税收；若生产区、船路或工匠网络不再合作，这种控制便会变得空洞。朱温、李克用、杨行密和王建的区域力量各有条件，不能用一个统一的“军阀经济”替代它们同物资链条的不同关系。

\subsection{唐亡前的礼仪仍在发挥作用}
唐末的礼仪并非在军人掌权后完全失去意义。册封、改元、祭祀、官服、诏令和受禅仪式，仍为不同集团提供能够声明关系的语言。一个地方政权继续使用唐年号，可能是在争取旧朝合法性、安抚官员或同邻近势力区别自己；接受后梁册命也可能是获得名位和贸易空间的策略。礼仪不能直接证明内心忠诚，却能影响谁被承认为官、谁可以发文书、谁有资格同外部政权交涉。

仪式的物质条件同样重要。诏书要有人抄写和递送，朝会需要场所、服饰和守卫，改元要通过地方文书进入实际日期。907年的受禅文本试图把军事胁迫包装为可读的秩序，正说明仪式仍有组织政治关系的力量；它无法消除军府对道路、城门和皇帝人身的控制。读者应同时看到这两面：礼仪不是虚假的装饰，也不是可以压倒军事实力的魔法。

唐号终结后，旧有礼仪和行政习惯会在不同区域继续被改写。五代十国中的官员、寺院、市场和地方家族并非从零开始生活；他们带着唐后期形成的文书、税制、道路与记忆进入新的政治世界。这种延续并不否认907年的暴力断裂，而是说明断裂的后果只能在随后许多年中逐渐显现。

\section{从军镇继承到海路来客：八个不在京城的晚唐场景}

\subsection{淄青的继承不是一封任命书能够解决的事}
建中二年，李正己死后，李纳在淄青平卢接过军府，\eventanchor{LT-781-LI-NA-QI}{李纳自立并与河北诸镇拒命（781）}。从长安看，这很容易被写成一个节度使之子不受朝命的瞬间；从淄青看，先要处理的却是军府里谁接管文书、仓库和军士，沿海与内陆的州县向谁交纳钱粮，父亲旧部愿意服从哪一种安排。李正己留下的并非一块可以完整继承的领地，而是一组由军队、地方官、乡里中介、运输路径和朝廷名义勉强接合的关系。新主人的名字写进军府公文，并不会自动让这些关系继续运转。

淄青平卢的地理使这种继承格外具体。它面对海滨，又连着山东、河南的道路；粮食、盐货、布帛、马匹和消息不必全经由长安才可流动。军府若要留住部众，必须使仓廪和赏给仍可兑现；地方人若要维持市场与耕作，也得判断哪一方的征发更迫近。唐廷要求李纳入朝或受命，表面上是在争一项官职，实则也在试图重新接通这些资源的出口。李纳同河北诸镇相互呼应，并不意味着他们共享一套稳固的政治计划；各镇担心的道路、粮源、邻镇和内部将校都不同。联盟首先是一种在压力下暂时互认的办法。

《旧唐书》《新唐书》的传记把李纳放入叛镇的谱系，《资治通鉴》则把同年的命令与对抗排成连续年序。这些文字让我们能够确定继承和拒命的大致次序，却多从中央失序的角度观看军镇。它们很少留下某一仓场怎样开支、一个县怎样分摊军粮、普通军户怎样看待继任者的记录。正因为如此，不能把“自立”译成一座城市在某日忽然脱离帝国，也不能把朝廷授官当作当地生活立即改写的证据。可见的是名义和军事资源开始脱节；不可见的部分，恰恰是这种脱节如何被无数小规模的交涉维持。

这一场继承还改变了后来朝廷谈判的尺度。朝廷若要恢复对河北的影响，不能只等待一位将领表示恭顺，而要面对军府已经能把地方征取、部众安置和区域交通连成自己的节奏。军镇也不能完全不使用唐的官号、任命和赦令，因为这些语言仍能帮助它约束部将、同邻镇交往并在地方社会取得认可。晚唐所谓中央与地方的对立，往往就发生在这种既相互借用又相互防备的关系中。李纳的继承不是唐朝衰落的一枚抽象注脚，而是一处可以看见权力怎样落在账簿、道路和人的去留上的现场。

\subsection{南诏遣使时，册封语言跨过的不是一条空白边界}
贞元十年，南诏同吐蕃关系破裂后遣使唐廷，唐朝以册封和使节往来重新安排西南关系，\eventanchor{LT-794-NANZHAO-TIBET-RUPTURE}{南诏与吐蕃关系破裂后遣使唐廷（794）}。使者抵达长安以前，已经穿过山地、河谷、关隘和一系列需要补给与翻译的节点。唐廷接到的不是一张天然指向自己的外交地图，而是一个在吐蕃、南诏、四川与安南诸方向之间重新衡量利害的政权。南诏选择遣使，意味着它认为改变关系可以带来军事、贸易或政治上的空间；唐廷选择回应，也是在西南压力之下寻找可用的伙伴，而不是单方面恢复了早已稳固的统属。

册封仪式很适合被后来的史书写成秩序恢复：皇帝给出称号，远方政权接受文字，双方似乎各归其位。但仪式之后，驿路是否安全、边城由谁供给、使团怎样往返、地方首领是否接受新的安排，仍要在具体地点逐一处理。四川通向云南的道路不只供使节行走，也牵动军粮、马匹、盐货和沿途社区。南诏与吐蕃关系的变动会使这些道路的风险重新分配；唐方在某处增加守备，未必能决定另一处山口的行动。把这一关系写成“南诏归附”，会把最需要解释的空间条件删去。

两《唐书》的南蛮传和《通鉴》保存了唐廷所看见的遣使、册封与同盟，它们的叙述语言自然把成功放在中原朝廷一边。立在南诏一侧的《南诏德化碑》则以王权自身的历史与功业组织记忆。两种材料并不需要被强行拼成一句毫无裂缝的话：它们恰好说明，同一外交行动在不同政治共同体中具有不同的名称和用途。碑文的王权修辞不能直接告诉我们所有边地居民的意见，唐廷的册命也不能直接证明南诏内部没有争论。读者应当保留这种差异，而不是让一方材料替另一方消失。

这次遣使把晚唐的边界问题带回了日常。边界不是一条由中央画出的线，而是商旅、使者、军人、翻译者和山地社区不断确认或绕开的通路。唐廷通过册封得到的是一个可以使用的外交接口，付出的则是承认西南局势无法只靠自己的军令处理。南诏得到的也不是一份永久安全，而是在多方压力中增加了一种可以谈判的关系。后来安南、四川和云南方向的冲突会继续暴露这种不稳定；794年的合作不能预告一切，却让我们看见战争与和平都要经过道路才能变成现实。

\subsection{成德的请罪，把长期战争压缩进几道文书}
元和十三年，王承宗遣使请罪，成德镇接受朝廷条件，\eventanchor{LT-818-CHENGDE-SUBMISSION}{王承宗遣使请罪，成德接受朝廷条件（818）}。在帝纪的句子里，遣使、请罪、赦宥和授官前后相连，仿佛一场长期的对抗在一纸文书抵达时便告结束。可是使者从成德到长安所带来的不只是措辞：它还带着军镇愿意在何种范围内承认朝廷、朝廷愿意在何种范围内保留地方既成安排的试探。多年战争之后，双方都要计算继续用兵的粮饷、道路和部众成本。请罪的语言提供了一个可以停手的形式，却不能抹掉这些成本已经怎样改变了地方。

成德所在的河北平原并不是任一方可以随意推移棋子的空地。州城、渡口、田地、军营和邻镇之间的距离，决定调兵或断粮并不只是发令问题。王承宗若要接受条件，必须让将校相信这不会危及他们的地位；朝廷若要赦宥，也得考虑如何使这一安排不被其他军镇视为软弱或威胁。所谓和解因此不是一方完全制服另一方，而是各自把难以承担的风险转为一种暂时可用的政治语言。对州县居民来说，文书上的和解可能减少某些战事，却不保证军粮、差役或旧日债务立刻消失。

《旧唐书》《新唐书》的王承宗传与《通鉴》能证实请罪和朝廷处置的主要次序，却都倾向把赦宥写成中央重新取得控制的节点。请罪、授官和赦书本来就是政治文本：它们告诉我们双方愿意公开承认什么，并不直接告诉我们成德境内的税收、军队和司法在次日怎样运行。地方文书、碑志和区域材料远不足以逐县补全这种变化，因此不能把沉默填成整齐的“复归”。承认这种限度，不会削弱事件的意义，反而能使读者看见一份赦令究竟能做到哪里。

元和时期的军事成功常被概括为唐廷中兴。成德的安排更适合被读作一种有代价的恢复能力：中枢仍可调动名义、官职、部分财政和军事压力，使一个强镇重新进入可谈判的秩序；它却没有获得无须协商的日常统治。这个制度收益是停止一段高成本冲突并重新打开命令渠道，副作用则是各方更清楚地知道官号、赦令和地方资源可以分开使用。几年后河北的军变会表明，纸面上的安定并未消除军营内部和地方财政的压力。读成德，不能只看结案，还要记住它留下的未决条件。

\subsection{魏博军营的夜晚，改变了河北的政治时间}
长庆二年，魏博、成德、卢龙相继发生军变，田布、田弘正等被杀或败亡，唐廷承认新镇将，\eventanchor{LT-822-HEBEI-MUTINIES}{河北诸镇军变与朝廷承认新镇将（822）}。当军士在营中决定杀害或驱逐上级时，朝廷此前给出的任命、赦书和道德评价忽然显得很远。营门、兵器、粮饷、同乡关系和谁能先取得城门，比诏书中的官名更直接地影响次日的局面。田弘正从魏博移镇成德，本想把中央任命带入一个长期自成节奏的军府；对留在魏博的军士而言，新的安排也可能意味着薪给、首领和保护关系发生不确定的移动。材料不允许我们替他们列出统一的要求，但军营的行动不能被简化成一种天生的叛逆。

军变扩散时，河北并没有变成一片没有秩序的黑暗。每一镇都要迅速决定谁发军令、谁守城、谁向周边传递消息、谁负责给部众筹粮。唐廷也不能只用“平乱”二字回应：远征需要钱粮，邻镇的立场会影响道路，贸然进入又可能把原本可谈的军府推向联合。承认新镇将看似退让，实际是一种在有限资源中避免更大断裂的选择。它让朝廷保留了名义上的关系和某些沟通渠道，却也承认自己无法规定所有军营的继承次序。

《旧唐书》穆宗本纪以及田弘正、王廷凑等传记，最清楚地保存了朝廷任命、军变和处置的顺序。它们把杀害上官的军士置于叛乱叙述中，这符合官修史书维护官僚与君臣秩序的需要。欠饷是否为普遍原因、地方税源怎样进入军营、战争怎样波及市镇与乡村，现有材料远没有同样清楚。因而不能用一项猜测解释全部军变，也不能从一镇的情形推出整个河北。史料的偏斜本身提示我们，军营中最能发声的是被承认为政治行动者的人，更多承担供给和逃难的人则只在后果中隐约出现。

这场军变把成德请罪后隐藏的矛盾重新翻到表面。朝廷可以令一位节度使入朝、移镇或受官，却未必能替他取得部众的信任；军镇可以拥立新帅，也仍需要唐的官号、邻镇的承认和地方的粮源。制度没有简单崩溃，而是变成多方都要反复验证的关系。对于后来的晚唐，河北的经验留下了一条冷峻的教训：若军队的供给、迁调和继承只在中枢文书上得到安排，军营会在自己的时间里作出答案，而那个答案往往先于朝廷抵达。

\subsection{圆仁从扬州北上，看见的不是一幅静止的盛唐图}
开成三年，日本僧圆仁随遣唐使抵达扬州，后入长安求法，\eventanchor{LT-838-ENNIN-ARRIVAL}{圆仁入唐求法与东亚交通观察（838）}。一名僧人下船时首先面对的不是“唐帝国”这个抽象名词，而是港口的停泊、行李的搬运、译语的困难、地方官的查验以及往北行进所需的凭验与食宿。圆仁的目的在于求法，他选择记下的寺院、师友、仪式和行程也因此带有明确的宗教视角；但正是这种有限视角，让一段常被中央本纪压缩掉的交通过程留下了纹理。海路并不是把日本和长安直接相连的蓝线，它以扬州为入口，再交给河道、道路、驿传和寺院网络逐段接续。

扬州的城市生活使这种接续可见。外来使团、商人、僧侣和地方居民使用的并非同一套资源：有人依靠官方接待，有人借寺院投宿，有人需要在市场购置物品或等待船只。城市能够容纳跨海来客，靠的是仓储、渡口、书写和熟悉规矩的人，而不是一种抽象的开放。圆仁所写的见闻不能代表所有来唐者，也不能让我们推断每一处港市都具有同样秩序；它仍提醒读者，唐后期的东亚联系在军镇、财政紧张与宫廷斗争之外持续存在。政治的不稳定会改变通行的成本，却不会自动使海陆网络停止。

从扬州到长安的路同样是一串不平等的门槛。寺院可以提供学习、介绍和住宿的关系，州县文书与官方礼仪则决定某些行动是否被承认。圆仁能以遣唐使成员和求法僧的身份移动，正说明身份会打开特定通道；没有留下日记的船户、翻译者、雇工和沿路居民，未必拥有同样的保护。不能把一位有文字能力的旅行者当作普通社会的代言人。不过，他留下的日期、地名与寺院活动，足以让我们问：消息和物品如何穿过不同权力区，宗教网络怎样同市场和官府相互借力。

《入唐求法巡礼行记》与《续日本后纪》的遣唐使记录从日本一侧保留了这段行程，《唐会要》关于诸蕃使的材料则显示唐廷如何分类来客。它们的视角不能彼此替换。唐廷的分类强调可管理的使节，圆仁的日记强调求法的困难与收获，日本朝廷的记载又服务于本国的遣使制度。把这些材料并置，读者会发现“文化交流”不是一种没有摩擦的美名，而是由航海风险、地方接待、宗教权威和政治许可共同组成的实践。晚唐的世界并不只向内收缩；它在一些港口、寺院与道路上，仍向外部保持着需要不断维护的开口。

\subsection{泽潞的继承危机，暴露了军府内部并不一致的利益}
会昌三年，昭义节度使刘从谏死后，养子刘稹拒命，唐廷发兵征讨泽潞，\eventanchor{LT-843-ZE-LU-CAMPAIGN}{刘稹拒命与唐廷征讨泽潞（843）}。泽潞的争端首先是一场继承问题：谁有资格接过军府，谁能让将校、幕僚和附近州县继续合作，谁能把先主留下的关系说成合法。朝廷以拒命命名刘稹，自有维护任命权的理由；军府内部却未必只在“顺”与“叛”之间选择。养子身份、旧部的依附、家族关系和军队供给，都可能使不同人对继承有不同判断。一个军镇的名字在地图上是完整的，它内部的利益却可能早已分岔。

泽、潞所在的山地与通道使战事不能只靠宣言结束。进军者要计算关隘、补给、马匹和季节，防守者也要维持城寨、传递消息并防止内部离散。唐廷能在会昌年间集中兵力，并不表示它可以无成本地把命令送进每一处山谷；地方将校若改变立场，也不必然是因为某一纸诏令突然具有魔力。军事压力、赦免承诺、利益分配和对前景的判断会共同影响选择。史书常把最后的结果组织成中央讨平叛镇的故事，读者则应在结果之前看见这些被道路和粮食限制的选择。

《旧唐书》武宗本纪、刘从谏传与《通鉴》提供了这场征讨的年序，却把刘稹的养子身份和“叛逆”身份紧密连在一起。这种写法能说明朝廷怎样为用兵辩护，不能直接说明泽潞的平民、商人或兵士如何理解继承，更无法可靠计算战争中的死亡与损失。地方独立材料的稀少是这里不能回避的空白。叙事不应因此虚构一座城中人人的情绪，也不应把少数将领的立场当作全部社会。能够确认的是，政治合法性在军府交接时会变成实际的供给与服从问题。

泽潞事件同会昌废佛常被放在同一位皇帝的“强势作为”之下，其实二者涉及不同的资源接口。前者考验朝廷能否在山地军镇的继承危机中集结兵力，后者考验它能否通过州县与寺院重新分类人口和资产。两种行动都显示武宗朝一度具备集中权力的能力，也都留下执行范围不均的疑问。将它们并列，不能得出一个全境恢复控制的结论；更合适的读法是，唐廷在若干节点重新取得主动，却仍须依赖地方中介、道路和可被调动的军队。强势从来不是均匀铺开的状态。

\subsection{浙东州县的危机，不应只从裘甫的名字开始和结束}
大中十二年，裘甫在浙东起事并攻扰州县，唐廷调兵镇压，\eventanchor{LT-858-QIU-FU-UPRISING}{裘甫起事与浙东州县危机（858）}。官修年表用首领姓名把这类事件固定下来，仿佛一群人只是从一个名字出发又在一个名字处消散。对于浙东的州县而言，危机首先表现为城门是否关闭、仓粮是否转移、渡口和山路能否通行、乡里如何应对征调与传言。当地的水网、丘陵、市场和海路让人群有多种移动的可能，也让追捕、供给和通讯都更加依赖熟悉地形的人。起事者能攻扰州县，并不意味着他们控制了整个区域；官军进驻，也不意味着地方关系立即恢复原状。

将裘甫及其追随者简单写成一种固定的社会群体尤其危险。正史材料没有给出足以辨认每个参与者出身、人数和动机的连续记录。有人可能因负担、债务、兵役或地方冲突加入，有人可能在军队抵达后离开或被迫协助；这些可能性不能被写成已知的集体宣言。可以较确定地说的是，州县秩序受到挑战后，地方官、军人、乡里首领和运输者都要重新选择如何分配粮食、守卫和信息。一个起事之所以能扩大，往往不是因为某个口号覆盖一切，而是因为这些日常接口在特定时刻失去协调。

《旧唐书》宣宗本纪与《通鉴》简要记下起事和朝廷应对，《新唐书》兵志的相关材料帮助我们把它放在军镇与调兵的背景中。它们让年序不至中断，却没有给我们一部浙东地方社会的逐日记录。海港遗址、钱币、窑业和地方文物能提醒我们该区域并非孤立，仍不能把某一处发现直接连到裘甫的队伍。史料不足要求的不是放弃叙述，而是缩小断言：我们写州县与道路受到扰动，不能编造每座村落的遭遇；我们写唐廷需要调兵，不能把平叛军数和伤亡数字说得确凿无疑。

裘甫起事也打破了把晚唐危机只放在河北或长安周边的习惯。浙东并非中央权力的安静背景，它有自己的生产、运输和地方组织，正因如此也会在征取、治安和政治竞争中承受不同的压力。一个区域经济活跃，不等于其居民免于军役和冲突；一场州县危机，也不自动证明整个江南已经失序。把这两种过度判断都放下，才能看见晚唐的多中心性：同一朝廷的年号之下，不同地区以不同速度感到军队、赋役和消息带来的震动。

\subsection{平叛之后的水路，仍要由地方人重新使用}
咸通元年，唐军平定裘甫余部，浙东起事在朝廷年表中告一段落，\eventanchor{LT-860-QIU-FU-SUPPRESSION}{唐军平定裘甫余部（860）}。一个“平”字很容易把读者带到终点，好像城门重新开启、官府恢复文书、农田与市场便回到冲突以前。实际上，军队离开或驻扎之后，最先需要恢复的是极其细碎的事情：渡口有没有船，仓里的粮是否还可用，商人敢不敢出行，逃离的人是否回来，地方官怎样重新催科。我们不知道这些问题在每一县得到的答案，也不能把恢复想象成整齐的庆典；正因无从逐一看见，更应当拒绝让“平定”掩去它们。

水路和沿岸市场在此时显出两面性。它们可以帮助官军运送人和粮，也让消息、逃亡者、货物与武装人员有更多绕行的机会。战时一处渡口的控制会改变相邻市镇的风险，战后同一渡口又要重新成为交易、课税和亲属往来的通道。地方社会不是被动等待军令的容器：船户、商人、乡里中介和寺院网络可能各自保存关于路线、债务和人情的知识。材料大多不替他们留下完整声音，却足以使我们不把道路当作地图上没有使用者的线。

《旧唐书》懿宗本纪和《通鉴》把平定写成朝廷取得的结果，叙述重心自然落在将领、命令和获胜上。它们不能用来量化死者、流离者或恢复所需的年月，也不能证明地方对官军具有一致的欢迎或敌意。后来的地方材料若记得某处修复，也只显示那一处的时间，不能自动延伸到整个浙东。来源的这种限制并非附带的谨慎语句，而会改变故事的重心：我们不把860年当作痛苦消失的日期，而把它当作新的协调工作开始的日期。

从浙东回看晚唐，中央并非只在长安才存在。它通过任命、调兵、税粮和年号进入州县，也在每次地方危机后暴露出自己必须借助当地道路、仓储和人际关系。地方也并非只在拒绝中央时才具有力量；它在恢复渡口、重开市场、安置军队和重新书写账册时持续塑造国家可以做到什么。裘甫事件的制度得失正在这里：平叛能够重新接通一段秩序，却不能免除维持这段秩序的日常成本。那些成本会在九世纪后半更大规模的兵变、饥荒与战争中，再次回到道路和城市之上。

\section{同一朝廷的年号下，边城、宫门与行在各有时间}

\subsection{幽州的边报，不能把回鹘余部写成一个整齐的敌人}
会昌三年，幽州节度使张仲武击退进入边境的回鹘余部，\eventanchor{LT-843-048}{张仲武击退回鹘余部（843）}。一封边报送到长安时，往往已把沿路遇见的不同人群、不同时间的冲突和不同地点的移动压缩为“破回鹘”几个字。幽州军府需要知道的却更细：哪一处关口可守，马匹从何处补充，商旅是否还敢通行，来者是寻求交易、避难，还是已进入武装冲突。回鹘汗国崩解后的北方流动不可能由一个名称概括；唐方文书中的类别首先服务于边防和朝廷分类，并不等于这些人的内部关系已经被我们看清。

幽州的力量来自它占据北方道路的能力，也受同一能力限制。守住一段路，需要军人、牧地信息、粮草和地方供应；一旦把大量资源投向边境，城内和周围州县的负担也会随之增加。张仲武的军功让朝廷可以把一次冲突写成秩序恢复，但军镇本身会借边防任务积累武装与谈判地位。中央需要它挡住风险，也要担心它因掌握道路而拥有过多自主空间。边镇并非只在对抗中央时才成为问题；它在履行中央最需要的职责时，同样会改变权力的分配。

《旧唐书》张仲武传、回鹘传与《通鉴》保存了这一年冲突的框架，但其中的俘获数和“破”字混合了多个部众、地点与行动。把报功数字直接换成精确战果，或把一切移动者都写作同一军事集团，都会超过材料的承载力。较可靠的读法是把843年看作北方道路重新紧张的一刻：唐廷、幽州军府和不同的草原人群都在调整安全、交换和驻守的安排。它告诉我们晚唐边境的变化，不是帝国地图边缘的一道涂色，而是道路上的人不断决定谁能通过、谁要付费、谁会被拒在门外。

\subsection{宣宗即位时，宫廷的继承也要经过人和门}
会昌六年三月，武宗去世，叔父李忱即位，是为宣宗，\eventanchor{LT-846-050}{武宗去世与宣宗即位（846）}。帝纪的换代往往写得极短：一位皇帝卒，一位新君立，年号和诏令继续向外发送。宫廷内部却不是一间会自行平稳运转的房屋。谁守宫门，谁保管符信，谁能接近皇族，谁把新诏传给外朝，都会影响继承能否被承认。宦官参与其中，并不意味着他们能够随意决定一切；皇室成员、文官、禁军和既有的礼仪程序同样构成限制。我们所能确证的是继承结果与大致时间，不能把后来的政治格局倒推成当夜每个人早已一致的意图。

武宗末年的宗教政策和军政举措仍在地方产生后续影响，新皇帝即位也不可能按一声号令把这些后果全部改写。地方官要重新理解哪些命令继续有效，寺院与家庭要面对已经发生的资产、人身和关系变化，边镇则会观察中枢是否改变任用、赏罚与供给的方式。因而继承对不同地点有不同的日期：长安宫门的消息先到，诏书经过驿路才成为州县可见的权威，乡里是否感到变化又取决于征收、司法和军役如何调整。把皇帝更替当作全国同时发生的转折，会失去这种传播过程。

《旧唐书》的武宗、宣宗本纪和《通鉴》为我们保留了宫廷更替的主线，细部则主要来自由宦官主导的继承叙事。它们能说明某些人怎样被放进合法的故事，不能完整揭示被排除的选择和沉默。来源限度并非让继承失去意义，反而使读者看到制度的实际用途：礼仪、诏令和官名把高度不透明的宫廷转折转换为外部可接受的连续性。唐廷可以借此避免权力真空，代价是更多实际决定仍要交给掌握门禁、文书和道路的人去落实。

\subsection{蔡州的军号，在黄巢危机中成为另一种资源凭据}
中和元年，蔡州节度使秦宗权据蔡州自立，\eventanchor{LT-881-QIN-ZONGQUAN}{秦宗权据蔡州自立（881）}。这一年黄巢军已在长安，僖宗行在成都；若只跟着两京的年序，蔡州的变化很容易被看作旁枝。对中原道路上的州县来说，它却意味着又一个能够征集兵粮、发布军令和争夺城镇的力量出现。军号本来是唐廷授予的行政与军事名义，在危机中也能被地方首领转化为组织部众、索取供给和同邻镇谈判的凭据。所谓自立，不是凭空创造权力，而是把既有的武装、仓储和关系从朝廷的框架中拔出，另作使用。

蔡州及其周边处在多路军队、逃亡者和粮道反复穿行的地带。一个军府要持续存在，不能只靠宣布不奉朝命；它还得让部众吃饭，控制或威胁邻近城镇，找到可以征取或交换的物资，并处理内部将校的忠诚。地方居民面对的选择往往没有史书中的旗帜那样清楚：留在城中可能遭遇征发，离开要承担道路风险，向一方提供信息又可能在另一方到来时成为罪状。我们没有材料逐户重建这些处境，因而不应把蔡州写成纯粹的暴力舞台；但也不能因细节不全就忽略军府扩张给日常安全带来的压力。

两《唐书》秦宗权传和《通鉴》保存了敌对军镇与后出朝廷叙事所强调的暴行。那些记载对于理解恐惧和战争语言重要，却不能把夸张的杀掠数字直接换算为人口损失统计，也不能把敌人笔下的形象当作唯一解释。更可把握的是，黄巢危机没有把地方力量压成两方对峙，反而使多个军府在中原争夺道路和资源。蔡州的出现使唐廷即便收复长安，也难以重新取得独立的中原控制；后来的军镇竞争正是在这些彼此叠加的战争空间里被推向新的规模。

\subsection{兴元行在的诏书，必须穿过关中各镇的许可}
光启元年，王重荣与田令孜冲突，僖宗离长安出奔兴元，\eventanchor{LT-885-067}{王重荣与田令孜冲突，僖宗出奔兴元（885）}。皇帝离开都城时，最先暴露的不是一个抽象的威望问题，而是队伍能否通过关中道路、沿途由谁供应、护卫是否可靠、消息如何不被截断。王重荣、田令孜以及周围军镇的冲突，把朝廷的财政、盐利、军权与私人关系缠在一起；史书常急于排列谁应承担责任，行在中的人首先要面对的却是下一处可停留的城、可用的粮与愿意放行的军队。兴元成为行在，并非因为它天然替代长安，而是因为当时的道路和地方力量让皇帝暂时能够在那里继续发令。

行在发出的诏书仍具有名义力量，但其力量要经由地方传递才会显现。某一节度使接受诏令，可能是在计算同皇帝的关系、同邻镇的竞争以及本镇的资源；另一镇拖延或拒绝，也未必只是公开的反叛，而可能是保留自身安全的策略。皇帝身边的宦官、文官、卫兵和随行供给者同样不是一个整体，他们的生存依赖于正在变化的关中网络。把僖宗出奔写成“朝廷逃亡”，会遮住行在仍在组织仪式、文书和任命；把它写成“权威依旧”，又会忽视每一道命令都须经过军府许可的脆弱。

《旧唐书》僖宗本纪、王重荣传和《通鉴》对冲突的责任归属具有强烈党派性，连出奔路线和各镇态度也须按时段分别核对。因此，我们不宜替田令孜、王重荣或每一镇预设唯一动机。事件的意义不在于为谁判定忠奸，而在于看见晚唐皇权如何被迫从固定都城变成移动的行政中心。制度上，行在使唐廷在失去长安时仍能维持连续名义；它的代价是宫廷越来越依赖能提供道路、营地和军粮的地方力量。这种依赖会在昭宗时期进一步加深。

\subsection{成都城门易手时，西川不是等待中原决定的空白地}
大顺二年，王建驱逐陈敬瑄，控制成都并成为西川实际主帅，\eventanchor{LT-891-070}{王建驱逐陈敬瑄，控制成都（891）}。成都的城门、府库、街市和通向蜀道的道路，在这场变化中都比一个后来王朝的名称更早具有意义。王建进入西川并非在空地上建立势力：陈敬瑄已有军政关系，地方官、士人、商人、军人和城外州县也已有各自的安排。谁能使守城者转向、谁能取得粮食和文书、谁能让周边承认新的命令，才决定“控制成都”在何种程度上成为现实。后来的前蜀史很容易把这些过程写成开国者应得的胜利，现场所面对的却是资源和关系的重新分配。

西川的地理为行在提供过屏障和供给，也使区域力量能够积累自己的节奏。山路并非天然的防线；它需要驿站、搬运、熟悉路线的人和可供停驻的城镇。成都的市场与仓储也不是只服务军队的器皿，它们连接手工业、农业、寺院和周围交通。军府接管城市时，新的守卫、征发和任命会进入这些关系，而关系本身也限制军府能做什么。把巴蜀说成富庶所以必然能自立，是另一种过度简化：资源必须被组织、运输和承认，战争与继承随时都可能使这种组织断裂。

《旧唐书》王建传、《新五代史》前蜀世家与《通鉴》保存了陈敬瑄退守和王建取得成都的年序，但主要成书于王建政权已经建立之后。它们对于成都居民的态度、地方妥协的细节和城市损失并不透明。我们可以确认权力中心发生转移，不能把它描成全体城民欢呼或全部秩序毁灭。正是这一份克制，让成都不再只是唐末帝纪的避难背景：它是一座有自身道路、市场和记忆的城市，其军府转换预告了十国时代区域政权将如何从晚唐已有的行政与经济网络中长出。

\subsection{太原城下的失败，使诏令的边界变得可见}
大顺元年，昭宗下诏讨伐李克用，朱温等军未能攻克太原，\eventanchor{LT-890-069}{昭宗讨李克用而未能攻克太原（890）}。一道讨伐诏书发出后，战争并不会沿着纸上的方向自动前进。通往河东的道路要让不同军镇的部队通过，军粮要在州县、仓场和山口之间接续，主将还要判断邻军是否真会在约定的时候抵达。太原周边的地形、城防和李克用在河东经营的关系，使进攻者不能只凭朝廷名义取得优势。诏令能够界定谁被称为“讨伐”的对象，却不能替军队运送一粒粮，也不能保证各镇把本地兵力交给同一作战计划。

对朱温以及其他奉诏出兵的军府而言，参与战争同时也是一场关于声望、补给和未来位置的计算。为朝廷效力可以取得名义和官职，长途进军却会消耗部众、马匹与地方财力；若同盟失败，谁承担退军和损失的责任又会成为新的冲突。李克用一方也不是只靠城墙等待结果，而要维持河东内部的军粮、将校与道路。晚唐军事集团之间的较量因此很少是“中央军”对“一名叛将”的简单对阵，而是多套供给与承诺在一个地理空间中互相试探。太原未被攻克，不只是一次战役的败报，也是各方无法把资源完全交给朝廷统一调度的证明。

《旧唐书》昭宗本纪、李克用传与《通鉴》保存了出兵和失利的基本年序，但兵力、战败责任和各镇实际投入多来自相互推诿的奏报与后来的叙述。不能据这些材料精确重建每支部队的规模，也不宜把失败归咎于某一位将领的胆怯或才能。较稳妥的结论是，昭宗仍能够用诏书召集一场跨军镇的行动，却不能保证这些被召集的力量成为一支稳定的共同军队。政治命令的可达范围与军事合作的可持续范围，已经明显不再重合。

太原城下的局面也让后来关中与中原的选择更容易理解。没有任何一个军府会只按“忠于唐”或“背叛唐”来安排自己：它们要观察道路、邻镇、部众和税源，并在朝廷的名义中寻找对自身有利的位置。唐廷通过出兵保留了裁决河东关系的主张，付出的代价则是让参与者更清楚地衡量中央无法独力承担战争。战役失败并没有立刻决定唐亡，却使下一阶段的联盟更短暂、军府更警惕，也使皇帝此后每一次调兵和迁徙都更依赖愿意合作的地方力量。这样的压力，最终会从太原周边的道路回流到长安、洛阳和各地的城门。
